\section{Estado del arte}
\label{sec:estado_arte}

La clasificación de textos financieros ha evolucionado desde sistemas basados en reglas hacia modelos de aprendizaje automático capaces de gestionar la ambigüedad del lenguaje natural.

\begin{itemize}
    \item \textbf{Técnicas de vectorización:} Para transformar texto en datos numéricos, destaca el uso de \textit{TF-IDF}, que asigna pesos a los términos según su relevancia en el corpus, y los \textit{Word Embeddings}, que permiten capturar relaciones semánticas complejas al representar palabras en espacios vectoriales continuos \cite{Minaee2021}.
    
    \item \textbf{Modelos de clasificación:} El uso de algoritmos clásicos como \textit{SVM} y \textit{Random Forest} sigue siendo fundamental en entornos que requieren eficiencia computacional y robustez en textos cortos. Estudios recientes de 2025 confirman que estos modelos ofrecen una mayor consistencia en la clasificación de etiquetas específicas frente a arquitecturas más pesadas \cite{John2025}. Aunque los \textit{Transformers} marcan el estado del arte en precisión, su alta demanda de recursos plantea desafíos de escalabilidad en aplicaciones ligeras \cite{Minaee2021}.
    
    \item \textbf{Aprendizaje incremental (Online Learning):} En el sector financiero, es crucial que los modelos se adapten a nuevos patrones de gasto sin reentrenamientos completos. La automatización de estos procesos requiere algoritmos optimizados (como SGD) que permitan un aprendizaje continuo y eficiente para garantizar la escalabilidad en sistemas de producción \cite{Oyewole2024, Kolluri2019}.
\end{itemize}